\documentclass[11pt]{article}
\usepackage[utf8]{inputenc}
\usepackage[T1]{fontenc}
\usepackage{lmodern}
\usepackage[margin=1in]{geometry}
\usepackage{amsmath,amssymb,amsthm}
\usepackage{enumitem}
\usepackage{booktabs}
\usepackage{xcolor}
\usepackage{tikz}
\usepackage[colorlinks=true,bookmarks=false,linkcolor=blue!50!black,urlcolor=blue!50!black]{hyperref}

\theoremstyle{plain}
\newtheorem{theorem}{Theorem}
\newtheorem{lemma}[theorem]{Lemma}
\newtheorem{claim}[theorem]{Claim}
\newtheorem{corollary}[theorem]{Corollary}
\theoremstyle{definition}
\newtheorem{definition}[theorem]{Definition}
\newtheorem{example}[theorem]{Example}
\theoremstyle{remark}
\newtheorem*{remark}{Remark}

\newcommand{\N}{\mathbb{N}}
\newcommand{\Z}{\mathbb{Z}}
\newcommand{\R}{\mathbb{R}}
\newcommand{\Q}{\mathbb{Q}}
\newcommand{\iv}[1]{[#1]}                 % Iverson bracket
\newcommand{\powerset}{\mathcal{P}}
\DeclareMathOperator{\ord}{ord}
\DeclareMathOperator{\gcdop}{gcd}

\setlength{\parindent}{0pt}
\setlength{\parskip}{0.6em}
\newcommand{\lesson}[2]{\href{../lessons/#1}{#2}}   % link to a lesson image

% Let long inline lists like \{1,2,5,6\} break across lines after a comma.
\mathchardef\mycomma="613B
\mathcode`,="8000
{\catcode`,=\active \gdef,{\mycomma\penalty0\hspace{0pt}}}
\emergencystretch=2em

\title{Group theory basics: orders, homomorphisms, isomorphism invariants}
\author{Notes from Carlos' lessons}
\date{}
\begin{document}
\maketitle

\textbf{Source lesson:} \lesson{2026-01-28-classification-of-groups.png}{2026-01-28 (classification of groups)}.

\section*{Definitions (from the board)}

A group $(G,\cdot)$ is associative, has an identity $e$ ($ge=eg=g$), and every element has an inverse ($g^{-1}g=gg^{-1}=e$). It is not necessarily commutative ($ab\ne ba$ may happen).
\begin{itemize}[nosep]
  \item \textbf{Order of the group} $|G|$: the number of elements.
  \item \textbf{Order of an element} $\ord(g)$: the least $k\ge1$ with $g^k=e$, or $\infty$ if no such $k$ exists.
  \item Small examples on the board: $\Z_1,\dots,\Z_4$ under $+$, with element orders $(1)$, $(1,2)$, $(1,3,3)$, $(1,4,2,4)$; and $(\R,+)$, where every non-zero element has order $\infty$.
\end{itemize}

\section*{Cancellation (the ``sudoku property'')}

\begin{claim}
$gh=gk\Rightarrow h=k$, and likewise on the right. So every row and column of the multiplication table of a finite group contains each element exactly once.
\end{claim}
\begin{proof}
Multiply on the left by $g^{-1}$.
\end{proof}

\section*{Every element of a finite group has finite order}

\begin{theorem}
If $G$ is finite, every $g\in G$ has $\ord(g)\le|G|$.
\end{theorem}
\begin{proof}
Consider $g^1,g^2,\dots,g^{|G|+1}$: that is $|G|+1$ elements of a set of size $|G|$, so by pigeonhole $g^i=g^j$ for some $1\le i<j\le|G|+1$. Cancelling $g^i$ gives $g^{j-i}=e$ with $1\le j-i\le|G|$.
\end{proof}
(The board restricts to $g^1,\dots,g^{|G|}$ and excludes the identity to reach the same conclusion; the version above is the cleanest.)

\section*{Cyclic behaviour}

For a fixed $g$ the powers $g^n$ either exhaust $G$ or they do not.
\begin{itemize}[nosep]
  \item \textbf{Finite $G$, powers exhaust $G$:} $g^n\leftrightarrow n$ matches $G$ with $(\Z_{|G|},+)$.
  \item \textbf{Infinite $G$, powers exhaust $G$:} $g^n\leftrightarrow n$ matches $G$ with $(\Z,+)$ (the board lists $\dots,g^{-2},g^{-1},g^0,g^1,g^2,\dots$).
\end{itemize}
The matching is a group isomorphism in the sense below, because $g^ig^j=g^{i+j}$.

\section*{Homomorphisms}

A \textbf{homomorphism} $\varphi:(G,\cdot)\to(H,*)$ is a map with $\varphi(ab)=\varphi(a)*\varphi(b)$ (it ``preserves structure''; ``relabelling'' on the board).

\begin{claim}[Consequences]
\begin{enumerate}[nosep]
  \item $\varphi(e_G)=e_H$.
  \item $\varphi(g^{-1})=\varphi(g)^{-1}$.
\end{enumerate}
\end{claim}
\begin{proof}
(1) $\varphi(e_G)=\varphi(e_Ge_G)=\varphi(e_G)\varphi(e_G)$; cancel $\varphi(e_G)$ in $H$.
(2) $\varphi(g)\varphi(g^{-1})=\varphi(gg^{-1})=\varphi(e_G)=e_H$.
\end{proof}

An \textbf{isomorphism} is a homomorphism with an inverse homomorphism ($\varphi^{-1}\circ\varphi=\mathrm{id}_G$); equivalently, a bijective homomorphism.

\section*{Invariants: orders}

\begin{claim}
If $\varphi:G\to H$ is an isomorphism, then $g$ and $\varphi(g)$ have the same order.
\end{claim}
\begin{proof}
$g^k=e_G\iff\varphi(g^k)=e_H\iff\varphi(g)^k=e_H$. The first ``$\Leftarrow$'' uses injectivity (only $e_G$ maps to $e_H$); the second equivalence is the homomorphism property.
\end{proof}

\textbf{Consequence (the lesson's strategy for classifying groups).} For a group $G$, compute the multiset $S(G)$ of the orders of its elements. Isomorphic groups have the same $S(G)$, so different multisets prove that two groups are \emph{not} isomorphic. For instance $S(\Z_3)=\{1,3,3\}\ne S(\Z_4)=\{1,4,2,4\}$.

\section*{Notes}
\begin{itemize}[nosep]
  \item The board is a definitions-and-first-lemmas lesson; the pigeonhole proof and the order invariant are the only real proofs, and both are given above.
  \item The board only claims that element orders are \emph{an} invariant under isomorphism; it does not claim that $S(G)$ determines the group.
\end{itemize}

\end{document}
